\documentclass[UTF8,11pt]{ctexart}
\usepackage[a4paper,margin=24mm]{geometry}
\usepackage{amsmath,amssymb,amsthm,mathtools,mathrsfs}
\usepackage[all]{xy}
\usepackage{xurl,hyperref,enumitem}
\usepackage{tikz}
\usetikzlibrary{arrows.meta,cd,decorations.markings,calc,patterns}
\hypersetup{hidelinks,breaklinks=true}
\setlength{\emergencystretch}{3em}
\allowdisplaybreaks
\newtheorem*{theorem}{Theorem}
\newtheorem*{thm}{Theorem}
\newtheorem*{lemma}{Lemma}
\newtheorem*{proposition}{Proposition}
\newtheorem*{prop}{Proposition}
\newtheorem*{corollary}{Corollary}
\newtheorem*{cor}{Corollary}
\newtheorem*{claim}{Claim}
\theoremstyle{definition}
\newtheorem*{definition}{Definition}
\newtheorem*{defn}{Definition}
\newtheorem*{remark}{Remark}
\newtheorem*{example}{Example}
\newcommand{\R}{\mathbb R}
\newcommand{\C}{\mathbb C}
\newcommand{\Z}{\mathbb Z}
\newcommand{\Q}{\mathbb Q}
\newcommand{\N}{\mathbb N}
\newcommand{\F}{\mathbb F}
\newcommand{\D}{\mathbb D}
\newcommand{\bB}{\mathbb B}
\newcommand{\bC}{\mathbb C}
\newcommand{\bR}{\mathbb R}
\newcommand{\bZ}{\mathbb Z}
\newcommand{\bN}{\mathbb N}
\newcommand{\bQ}{\mathbb Q}
\newcommand{\bD}{\mathbb D}
\newcommand{\bF}{\mathbb F}
\newcommand{\CP}{\mathbb{CP}}
\newcommand{\RP}{\mathbb{RP}}
\newcommand{\sO}{\mathscr O}
\newcommand{\sI}{\mathscr I}
\newcommand{\sF}{\mathscr F}
\newcommand{\sG}{\mathscr G}
\newcommand{\sH}{\mathscr H}
\newcommand{\sR}{\mathscr R}
\newcommand{\sM}{\mathscr M}
\newcommand{\sV}{\mathscr V}
\newcommand{\sU}{\mathscr U}
\DeclareMathOperator{\Hom}{Hom}
\DeclareMathOperator{\Ext}{Ext}
\DeclareMathOperator{\Tor}{Tor}
\DeclareMathOperator{\Spec}{Spec}
\DeclareMathOperator{\Proj}{Proj}
\DeclareMathOperator{\supp}{supp}
\DeclareMathOperator{\im}{im}
\DeclareMathOperator{\id}{id}
\DeclareMathOperator{\Tr}{Tr}
\DeclareMathOperator{\tr}{tr}
\DeclareMathOperator{\rank}{rank}
\DeclareMathOperator{\ord}{ord}
\DeclareMathOperator{\dist}{dist}
\DeclareMathOperator{\End}{End}
\DeclareMathOperator{\Aut}{Aut}
\DeclareMathOperator{\Gal}{Gal}
\DeclareMathOperator{\vol}{vol}
\DeclareMathOperator{\Cl}{Cl}
\DeclareMathOperator{\coker}{coker}
\DeclareMathOperator{\codim}{codim}
\DeclareMathOperator{\divergence}{div}
\DeclareMathOperator{\Ric}{Ric}
\DeclareMathOperator{\grad}{grad}
\DeclareMathOperator{\Hess}{Hess}
\newcommand{\abs}[1]{\left\lvert#1\right\rvert}
\newcommand{\sabs}[1]{\lvert#1\rvert}
\newcommand{\babs}[1]{\bigl\lvert#1\bigr\rvert}
\newcommand{\Babs}[1]{\Bigl\lvert#1\Bigr\rvert}
\newcommand{\bbabs}[1]{\biggl\lvert#1\biggr\rvert}
\newcommand{\BBabs}[1]{\Biggl\lvert#1\Biggr\rvert}
\newcommand{\norm}[1]{\left\lVert#1\right\rVert}
\newcommand{\snorm}[1]{\lVert#1\rVert}
\newcommand{\bnorm}[1]{\bigl\lVert#1\bigr\rVert}
\newcommand{\Bnorm}[1]{\Bigl\lVert#1\Bigr\rVert}
\newcommand{\bbnorm}[1]{\biggl\lVert#1\biggr\rVert}
\newcommand{\BBnorm}[1]{\Biggl\lVert#1\Biggr\rVert}
\newcommand{\widebar}[1]{\overline{#1}}
\newcommand{\nicefrac}[2]{\frac{#1}{#2}}
\newcommand{\linnprod}[2]{\langle#1,#2\rangle}
\newcommand{\myindex}[1]{#1}
\newcommand{\glsadd}[1]{}
\newcommand{\avoidbreak}{}
\newcommand{\sourceref}[1]{\textnormal{[原文：\nolinkurl{#1}]}}
\newcommand{\thmref}[1]{\sourceref{#1}}
\newcommand{\lemmaref}[1]{\sourceref{#1}}
\newcommand{\propref}[1]{\sourceref{#1}}
\newcommand{\corref}[1]{\sourceref{#1}}
\newcommand{\defnref}[1]{\sourceref{#1}}
\newcommand{\exampleref}[1]{\sourceref{#1}}
\newcommand{\exerciseref}[1]{\sourceref{#1}}
\newcommand{\figureref}[1]{\sourceref{#1}}
\newcommand{\sectionref}[1]{\sourceref{#1}}
\newcommand{\appendixref}[1]{\sourceref{#1}}
\newcommand{\stacksref}[2]{\href{https://stacks.math.columbia.edu/tag/#1}{#2 [#1]}}


\newcommand{\E}{\mathbb E}
\newcommand{\Prob}{\mathbb P}
\newcommand{\Reff}{\mathcal R_{\mathrm{eff}}}
\newcommand{\eps}{\varepsilon}
\DeclareMathOperator{\diam}{diam}
\DeclareMathOperator{\Var}{Var}
\DeclareMathOperator{\Crit}{Crit}
\newcommand{\dd}{\,\mathrm d}


\begin{document}
\section*{078\quad 满足可证性条件的算术理论的Lob反射定理}
\subsection*{来源与整理范围}
Open Logic Project，\emph{The Open Logic Text}，Incompleteness / incompleteness-provability，lob-thm.tex（Git blob 5b93347f4d9d19573aa8e64844edeb881bb3dba3），附fixed-point-lemma.tex的作者证明及provability-conditions.tex所列P1--P3。
\par\url{https://github.com/OpenLogicProject/OpenLogic/blob/master/content/incompleteness/incompleteness-provability/lob-thm.tex}
\par\textbf{恢复方式（整理者说明）：}依据人类原文重新整理以补齐缺失题号；并非旧题证文件的逐字节恢复；2026-09-28。
\subsection*{作者命题与人类证明}

\paragraph{Notation (editor's expansion of the source macros).}
$\mathrm Q$ denotes Robinson arithmetic. For a formula $A$, $\ulcorner A\urcorner$ is the standard numeral of its G\"odel number; $\operatorname{diag}$ is the primitive recursive diagonal-substitution function. The formula $D_{\operatorname{diag}}(x,y)$ represents this function in $\mathrm Q$.
\begin{lemma}[Fixed-point lemma; the author's support proof]
Let $B(x)$ be any formula with one free variable $x$. Then there is a sentence $A$ such that $\mathrm Q\vdash A\leftrightarrow B(\ulcorner A\urcorner)$.
\end{lemma}
\begin{proof}
Given $B(x)$, let $E(x)$ be the formula $\exists y(D_{\operatorname{diag}}(x,y)\land B(y))$ and let $A$ be its diagonalization, i.e., the formula $E(\ulcorner E(x)\urcorner)$.
Since $D_{\operatorname{diag}}$ represents $\operatorname{diag}$, and $\operatorname{diag}(\ulcorner E(x)\urcorner)=\ulcorner A\urcorner$, $\mathrm Q$ can derive
\begin{align}
&D_{\operatorname{diag}}(\ulcorner E(x)\urcorner,\ulcorner A\urcorner),\label{repdiag1}\\
&\forall y(D_{\operatorname{diag}}(\ulcorner E(x)\urcorner,y)\to y=\ulcorner A\urcorner).\label{repdiag2}
\end{align}
Now we show that $\mathrm Q\vdash A\leftrightarrow B(\ulcorner A\urcorner)$. We argue informally, using just logic and facts derivable in $\mathrm Q$.
First, suppose $A$, i.e., $E(\ulcorner E(x)\urcorner)$. Going back to the definition of $E(x)$, we see that $E(\ulcorner E(x)\urcorner)$ just is
\[\exists y(D_{\operatorname{diag}}(\ulcorner E(x)\urcorner,y)\land B(y)).\]
Consider such a $y$. Since $D_{\operatorname{diag}}(\ulcorner E(x)\urcorner,y)$, by \eqref{repdiag2}, $y=\ulcorner A\urcorner$. So, from $B(y)$ we have $B(\ulcorner A\urcorner)$.
Now suppose $B(\ulcorner A\urcorner)$. By \eqref{repdiag1}, we have
\[D_{\operatorname{diag}}(\ulcorner E(x)\urcorner,\ulcorner A\urcorner)\land B(\ulcorner A\urcorner).\]
It follows that
\[\exists y(D_{\operatorname{diag}}(\ulcorner E(x)\urcorner,y)\land B(y)).\]
But that's just $E(\ulcorner E(x)\urcorner)$, i.e., $A$.
\end{proof}

\paragraph{Derivability conditions (assumptions, from the source).}
For the provability formula $\operatorname{Prov}_T$, assume:
\begin{enumerate}[label=P\arabic*.]
\item If $T\vdash A$, then $T\vdash\operatorname{Prov}_T(\ulcorner A\urcorner)$.
\item $T\vdash\operatorname{Prov}_T(\ulcorner A\to B\urcorner)\to(\operatorname{Prov}_T(\ulcorner A\urcorner)\to\operatorname{Prov}_T(\ulcorner B\urcorner))$.
\item $T\vdash\operatorname{Prov}_T(\ulcorner A\urcorner)\to\operatorname{Prov}_T(\ulcorner\operatorname{Prov}_T(\ulcorner A\urcorner)\urcorner)$.
\end{enumerate}
\begin{theorem}[L\"ob's theorem; the counted problem]
Let $T$ be an axiomatizable theory extending $\mathrm Q$, and suppose $\operatorname{Prov}_T(y)$ is a formula satisfying conditions P1--P3. If $T$ derives $\operatorname{Prov}_T(\ulcorner A\urcorner)\to A$, then in fact $T$ derives $A$.
\end{theorem}
\begin{proof}
Suppose $A$ is a sentence such that $T$ derives $\operatorname{Prov}_T(\ulcorner A\urcorner)\to A$. Let $B(y)$ be the formula $\operatorname{Prov}_T(y)\to A$, and use the fixed-point lemma to find a sentence $D$ such that $T$ derives $D\leftrightarrow B(\ulcorner D\urcorner)$. Then each of the following is derivable in $T$:
\begin{align}
&D\leftrightarrow(\operatorname{Prov}_T(\ulcorner D\urcorner)\to A)\label{L1}\\
&\qquad\text{$D$ is a fixed point of $B(y)$}\notag\\
&D\to(\operatorname{Prov}_T(\ulcorner D\urcorner)\to A)\label{L2}\\
&\qquad\text{from \eqref{L1}}\notag\\
&\operatorname{Prov}_T(\ulcorner D\to(\operatorname{Prov}_T(\ulcorner D\urcorner)\to A)\urcorner)\label{L3}\\
&\qquad\text{from \eqref{L2} by condition P1}\notag\\
&\operatorname{Prov}_T(\ulcorner D\urcorner)\to\operatorname{Prov}_T(\ulcorner\operatorname{Prov}_T(\ulcorner D\urcorner)\to A\urcorner)\label{L4}\\
&\qquad\text{from \eqref{L3} using condition P2}\notag\\
&\operatorname{Prov}_T(\ulcorner D\urcorner)\to\bigl(\operatorname{Prov}_T(\ulcorner\operatorname{Prov}_T(\ulcorner D\urcorner)\urcorner)\to\operatorname{Prov}_T(\ulcorner A\urcorner)\bigr)\label{L5}\\
&\qquad\text{from \eqref{L4} using P2 again}\notag\\
&\operatorname{Prov}_T(\ulcorner D\urcorner)\to\operatorname{Prov}_T(\ulcorner\operatorname{Prov}_T(\ulcorner D\urcorner)\urcorner)\label{L6}\\
&\qquad\text{by derivability condition P3}\notag\\
&\operatorname{Prov}_T(\ulcorner D\urcorner)\to\operatorname{Prov}_T(\ulcorner A\urcorner)\label{L7}\\
&\qquad\text{from \eqref{L5} and \eqref{L6}}\notag\\
&\operatorname{Prov}_T(\ulcorner A\urcorner)\to A\label{L8}\\
&\qquad\text{by assumption of the theorem}\notag\\
&\operatorname{Prov}_T(\ulcorner D\urcorner)\to A\label{L9}\\
&\qquad\text{from \eqref{L7} and \eqref{L8}}\notag\\
&(\operatorname{Prov}_T(\ulcorner D\urcorner)\to A)\to D\label{L10}\\
&\qquad\text{from \eqref{L1}}\notag\\
&D\label{L11}\\
&\qquad\text{from \eqref{L9} and \eqref{L10}}\notag\\
&\operatorname{Prov}_T(\ulcorner D\urcorner)\label{L12}\\
&\qquad\text{from \eqref{L11} by condition P1}\notag\\
&A\qquad\text{from \eqref{L9} and \eqref{L12}.}\notag
\end{align}
\end{proof}

\subsection*{整理说明与先修范围（非作者正文）}
按此前取得的作者题证恢复整理。保留Lob证明中的全部可证性推导，附不动点构造；P1--P3是本题明示假设，未声称这里证明了某个具体算术化满足这些条件。原文最后一步引用L-8与L-12，整理时将引用改为前文实际已写出的L-9与L-12，此为引用号校正，不增加证明步骤。宏展开、中文来源与记号段为编辑处理。标准先修为一阶算术可表示性和经典逻辑；难点是在对象语言内部协调对角句、可证性及两层引用，而非直接引用反射结论。
\subsection*{可修改的简短标注（非作者正文）}
$c$：算术理论内部可证反射原则的刻画。

$a$：语法算术化；对角不动点；Hilbert--Bernays可证性条件；内部化演绎。

\texttt{cross\_theory\_status = yes}。把以元理论形式给定的推导性质转为理论内部的可证性蕴涵，选取与目标A相匹配的不动点句，最后由该句的可证性返回A。位置：L1--L12。

全部标注为非穷尽、可修改初稿。
\subsection*{来源许可与编辑归属}
Open Logic Project，\emph{The Open Logic Text}，CC BY 4.0；许可见仓库README：\url{https://github.com/OpenLogicProject/OpenLogic/blob/master/README.md}。正文来源宏展开为标准LaTeX；编辑题号、取材范围与标注另列。
ChatGPT 加入题号、中文来源说明、整理范围和初稿标注；编辑说明与人类证明分列。
\end{document}
